\chapter{Equity Strategies}

This chapter formalises the equity-driven strategies implemented in \texttt{src/strategies}. Each subsection states the signal equations, execution logic, and representative parameter choices drawn from the production configuration.

\section{Cointegrated Pairs Mean Reversion}
The \texttt{PairsTradeStrategy} maintains a market-neutral spread between equities $S^A_t$ and $S^B_t$ with hedge ratio $\hat{\beta}$ computed offline. The spread process is
\begin{equation}
X_t = S^A_t - \hat{\beta} S^B_t.
\end{equation}
Assuming $(S^A_t, S^B_t)$ are cointegrated, $X_t$ is stationary with mean $\mu$ and standard deviation $\sigma$. Online statistics are re-estimated over a rolling window $\mathcal{W}$ via
\begin{equation}
\hat{\mu}_t = \frac{1}{|\mathcal{W}|} \sum_{u \in \mathcal{W}} X_u, \quad
\hat{\sigma}_t = \sqrt{\frac{1}{|\mathcal{W}|} \sum_{u \in \mathcal{W}} (X_u - \hat{\mu}_t)^2}.
\end{equation}
A z-score normalisation triggers trades:
\begin{equation}
Z_t = \frac{X_t - \hat{\mu}_t}{\hat{\sigma}_t}.
\end{equation}
With entry threshold $\theta_{\text{in}}$ and exit threshold $\theta_{\text{out}}$, orders follow
\begin{align*}
Z_t \geq \theta_{\text{in}} &\Rightarrow \text{short spread (sell $S^A$, buy $\hat{\beta}$ shares of $S^B$)},\\
Z_t \leq -\theta_{\text{in}} &\Rightarrow \text{long spread (buy $S^A$, sell $\hat{\beta}$ shares of $S^B$)},\\
|Z_t| \leq \theta_{\text{out}} &\Rightarrow \text{close outstanding position}.
\end{align*}

\paragraph{Example.} For symbols ``GS'' and ``MS'' with $\hat{\beta}=0.84$, $\theta_{\text{in}}=1.25$, $\theta_{\text{out}}=0.25$, and position size $q=10$, observing $Z_t=1.40$ generates orders $(-q, +q\hat{\beta}) = (-10, +8)$.

\paragraph{Implementation.} Listing~\ref{lst:pairs_zscore} shows the core computation in \texttt{src/strategies/pairs\_trade.py}.
\begin{lstlisting}[language=Python,firstnumber=340,label={lst:pairs_zscore},caption={Z-score evaluation in PairsTradeStrategy}]
spread = price_a - self.hedge_ratio * price_b
z_score = (spread - self.spread_mean) / self.spread_std
signals = self._evaluate_rules(z_score)
\end{lstlisting}

\section{Donchian Channel Breakout}
The \texttt{DonchianBreakoutStrategy} trades momentum breakouts using the Donchian envelope defined by the rolling extrema of closing prices $C_t$ over lookback $n$:
\begin{equation}
U_t = \max_{u \in [t-n, t-1]} C_u, \qquad L_t = \min_{u \in [t-n, t-1]} C_u.
\end{equation}
Signals are
\begin{align*}
C_t > U_t &\Rightarrow \text{enter long},\\
C_t < L_t &\Rightarrow \text{enter short},\\
\text{otherwise} &\Rightarrow \text{hold position}.
\end{align*}
Positions persist until a counter-signal occurs.

\paragraph{Example.} With $n=20$ and quantity $q=10$, if $C_t$ closes 1.2\% above $U_t$, the strategy buys $q$ shares; a drop below $L_t$ later flips to a short.

\paragraph{Implementation.} The vectorised backtest (Listing~\ref{lst:donchian}) multiplies the signalled stance with forward returns to obtain trade-level performance for validation routines.
\begin{lstlisting}[language=Python,firstnumber=87,label={lst:donchian},caption={Vectorised Donchian backtest kernel}]
upper = df['close'].rolling(lookback).max().shift(1)
lower = df['close'].rolling(lookback).min().shift(1)
signals[df['close'] > upper] = 1
signals[df['close'] < lower] = -1
strategy_returns = signals * market_returns
\end{lstlisting}

\section{Macro Lead--Lag Arbitrage}
The \texttt{MacroArbitrageStrategy} exploits delayed transmission from a leader asset $L$ to a laggard $G$. Over a lookback window of $m$ minutes the leader return is
\begin{equation}
R^{L}_{t,m} = \frac{P^L_t}{P^L_{t-m}} - 1,
\end{equation}
while the laggard move is constrained by $|R^{G}_{t,m}| < \delta$ to ensure it has not reacted. A z-score computed on the leader--laggard spread,
\begin{equation}
Z_t = \frac{R^{L}_{t,m} - \E[R^{L}_{t,m}]}{\sqrt{\Var(R^{L}_{t,m})}},
\end{equation}
triggers trades when $|Z_t| > \theta$. Qualifying signals are filtered by natural language sentiment scores $\sigma_t$ derived from partner news feeds to suppress idiosyncratic jumps.

\paragraph{Execution.} Upon a positive shock in the leader ($Z_t>\theta$), the bot shorts the leader and buys the laggard, sized to equal dollar exposure. Positions are held for $h$ minutes (default $h=120$) or until the spread normalises.

\section{Sector Momentum Rotation}
The \texttt{SectorMomentumStrategy} implements monthly relative-strength rotation among the SPDR sector universe $\mathcal{U}$. For each candidate $s \in \mathcal{U}$ the six-month momentum score is
\begin{equation}
M_s = \frac{C_s^{(t)}}{C_s^{(t-126)}} - 1.
\end{equation}
The top $N$ sectors by $M_s$ are retained, subject to a regime filter that requires the closing price to exceed its 200-day simple moving average ($C^{(t)}_s > \text{SMA}_{200,s}(t)$). Assets failing the filter are replaced with the defensive Treasury ETF $D$.

\paragraph{Allocation.} Capital $A$ is split equally across the final basket $\mathcal{B}$, yielding target weights $w_s = 1/|\mathcal{B}|$. The rebalance transacts $\Delta q_s = (w_s A / C_s^{(t)}) - q_s^{\text{prev}}$ shares.

\paragraph{Example.} If the top momentum names are XLK, XLY, and XLF, but XLY trades below its SMA$_{200}$, the portfolio holds \{XLK, SHV, XLF\} with one-third allocation each.
